\documentclass[11pt, twocolumn, landscape, a4paper]{article}

\usepackage[margin=1cm]{geometry}
\usepackage[utf8]{inputenc}
\usepackage[spanish]{babel}
\usepackage{amssymb, amsmath, amsbsy}
\usepackage{tasks}
\usepackage{enumerate}
\usepackage{graphicx}

\usepackage[pdftex]{hyperref}
\hypersetup{pdftitle={Filosofía y Lógica},pdfauthor={Luis Carrillo Gutiérrez}}

\fontfamily{cmss}\selectfont
\renewcommand{\familydefault}{\sfdefault}
\pagestyle{empty}

\begin{document}
\subsection*{Filosofía y Lógica - Repaso}


\begin{itemize}
\item{El gobierno ordena el servicio militar obligatorio para todos los jóvenes. Esta es una función del lenguaje :
\begin{tasks}(3)
\task{Informativo}
\task{Expresivo}
\task{Lógico}
\task{Directivo}
\task{Científico}
\end{tasks}
}
\item{La forma ``\textrm{P}'' :
\begin{tasks}(2)
\task{Es una proposición verdadera}
\task{Es una proposición afirmativa}
\task{Es una proposición falsa}
\task{Es una proposición simple}
\task{Es solo un símbolo proposicional}
\end{tasks}
}
\item{Constituye parte importante del ``Órganon'' y caracteriza a la lógica Aristotélica.
\begin{tasks}(2)
\task{Silogismo}
\task{Intuición}
\task{Inducción}
\task{Simbolismo}
\task{Proposición}
\end{tasks}
}
\item{Busco convertir (a la filosofía) en una ciencia similar a la matemática, relacionando para ello un lenguaje especial, fue :
\begin{tasks}(3)
\task{Sócrates}
\task{Aristóteles}
\task{Sto. Tomas}
\task{Lukasiewcz}
\task{Leibniz}
\end{tasks}
}
\item{De las siguientes definiciones sobre lógica, una es verdadera :
\begin{tasks}(1)
\task{Ciencia del cálculo racional cuantitativo.}
\task{Estudia la facultad del razonamiento.}
\task{Analiza el pensamiento objetivo.}
\task{Estudia la necesidad del razonamiento.}
\task{Ciencia de los métodos y leyes del razonamiento.}
\end{tasks}
}
\item{La formulación del cuadro de oposiciones que muestra las relaciones entre proposiciones categóricas, fue planteada por :
\begin{tasks}(5)
\task{Porfirio}
\task{Leibniz}
\task{Sócrates}
\task{Boeccio}
\task{Morgan}
\end{tasks}
}
\item{Se sabe que :
\begin{enumerate}[i.]
\item{José trabaja en el taller ya que no estudia en el colegio.}
\item{José estudia en el colegio sí y sólo si trabaja en el taller.}
\end{enumerate}
\begin{tasks}(1)
\task{José no estudia en el colegio.}
\task{José estudia en el colegio o trabaja en el taller.}
\task{José no trabaja en el taller.}
\task{José no estudia en el colegio cuando trabaja en el taller.}
\task{José estudia en el colegio y trabaja en el taller.}
\end{tasks}
}
\item{La proposición es un segmento lingûístico :
\begin{tasks}(1)
\task{Que carece de significado; pero es V ó F.}
\task{Que expresa orden o interrogación.}
\task{Que carece de significado y no es ni V ó F.}
\task{Que contiene símbolos que no necesita.}
\task{Es portador de significado y es V ó F.}
\end{tasks}
}
\end{itemize}
\end{document}